\documentclass[12pt,a4paper]{article}

\usepackage[brazilian]{babel}

\usepackage{amsmath}
\usepackage{amssymb}
\usepackage{enumitem}
\usepackage{etoolbox}
\usepackage{tocloft}
\usepackage{tikz}
\usepackage[most]{tcolorbox}
\usepackage{xcolor}
\usetikzlibrary{decorations.pathmorphing}

\setlength{\cftsecnumwidth}{3.2cm}
\setlength{\cftsubsecnumwidth}{1.5cm}

\renewcommand{\thesection}{Capítulo \arabic{section} -}
\renewcommand{\thesubsection}{\arabic{section}.\arabic{subsection}}
\pretocmd{\subsection}{\clearpage}{}{}

\AtBeginEnvironment{enumerate}{\pretocmd{\item}{\setcounter{equation}{0}}{}{}}

% Marcador Q.E.D. / Fim de Resolução em equações em display
\newcommand{\qedsymbolmath}{\tag*{\textcolor{blue!70!black}{\rule{1.6ex}{1.6ex}}}}

% Caixa para ferramentas e recursos teóricos
\newtcolorbox{recursosbox}[1][]{
    enhanced,
    colback=blue!3!white,
    colframe=blue!60!black,
    coltitle=white,
    fonttitle=\bfseries,
    title=Ferramentas e Recursos da Seção,
    arc=2mm,
    boxrule=0.8pt,
    left=6pt, right=6pt, top=6pt, bottom=6pt,
    #1
}

% Caixa para alertas de erros frequentes e dicas práticas
\newtcolorbox{dicasbox}[1][]{
    enhanced,
    colback=orange!4!white,
    colframe=orange!80!black,
    coltitle=white,
    fonttitle=\bfseries,
    title=Atenção aos Detalhes e Dicas de Resolução,
    arc=2mm,
    boxrule=0.8pt,
    left=6pt, right=6pt, top=6pt, bottom=6pt,
    #1
}

\title{
    \Huge \textbf{Um Curso de Cálculo Volume 1} \\
    \vspace{0.3cm}
    \Large 6\textordfeminine{} Edição (Guidorizzi) \\
    \vspace{0.5cm}
    \large \textit{Resoluções Completas e Passo a Passo do Capítulo 7 Derivadas}
}
\author{Rogério Campos}

\begin{document}

\maketitle
\newpage

\tableofcontents
\newpage

\setcounter{section}{6}
\setcounter{subsection}{0}

\section{Derivadas}

\subsection{Introdução}
Esta subseção não possui exercícios

\subsection{Derivada de uma função}

\begin{recursosbox}
\textbf{Definição Formal de Derivada num Ponto:}
Sejam $f$ uma função e $p$ um ponto de seu domínio. O limite
\[
\lim_{x \to p} \frac{f(x) - f(p)}{x - p}
\]
quando existe e é finito, denomina-se \textbf{derivada de $f$ em $p$} e indica-se por $f'(p)$. Se $f$ admite derivada em $p$, dizemos que $f$ é \textbf{derivável} (ou diferenciável) em $p$.

\vspace{0.2cm}
\textbf{Formas Equivalentes de Cálculo:}
Fazendo a mudança de variável $x = p + h$ (onde $x \to p \iff h \to 0$), obtém-se a forma alternativa clássica por acréscimo:
\begin{itemize}[noitemsep, topsep=2pt]
    \item \textbf{Forma por aproximação pontual:}
    \[
    f'(p) = \lim_{x \to p} \frac{f(x) - f(p)}{x - p}
    \]
    \item \textbf{Forma por acréscimo ($h = \Delta x$):}
    \[
    f'(p) = \lim_{h \to 0} \frac{f(p + h) - f(p)}{h}
    \]
\end{itemize}

\vspace{0.2cm}
\textbf{Interpretação Geométrica e Reta Tangente:}
\begin{itemize}[noitemsep, topsep=2pt]
    \item $f'(p)$ representa geometricamente o \textbf{coeficiente angular} ($m$) da reta tangente ao gráfico de $f$ no ponto $(p, f(p))$.
    \item A equação da reta tangente ao gráfico de $f$ em $(p, f(p))$ é dada por:
    \[
    y - f(p) = f'(p)(x - p)
    \]
\end{itemize}
\end{recursosbox}

\begin{dicasbox}
\textbf{Dificuldades Frequentes e Estratégias de Resolução:}
\begin{itemize}[noitemsep, topsep=2pt]
    \item \textbf{Escolha da forma do limite:} Para pontos numéricos fixos (ex.: $p = 2$), a forma $\frac{f(x) - f(p)}{x - p}$ frequentemente simplifica fatorações polinomiais mais rápido. Já para encontrar a função derivada geral $f'(x)$, a forma por acréscimo com $h \to 0$ costuma ser mais direta.
    \item \textbf{Cuidado com o quociente indeterminante:} O limite da derivada sempre produz a forma $\left[\frac{0}{0}\right]$. Você \textbf{deve} manipular algebricamente a expressão (fatorando, racionalizando ou cancelando termos) até eliminar o fator $(x - p)$ ou o denominador $h$ antes de avaliar o limite.
    \item \textbf{Funções por várias sentenças ou módulos:} Para verificar derivabilidade em pontos de transição (ex.: em $x = 0$ para $f(x) = |x|$), é obrigatório calcular os \textbf{limites laterais} do quociente de Newton:
    \[
    f'_+(p) = \lim_{x \to p^+} \frac{f(x) - f(p)}{x - p} \quad \text{e} \quad f'_-(p) = \lim_{x \to p^-} \frac{f(x) - f(p)}{x - p}
    \]
    $f'(p)$ existirá se, e somente se, $f'_+(p) = f'_-(p)$ e ambos forem finitos.
    \item \textbf{Derivabilidade implica Continuidade:} Se $f$ é derivável em $p$, então $f$ é automaticamente contínua em $p$. Se a função apresentar descontinuidade em $p$, a derivada em $p$ não existirá imediatamente (sem necessidade de calcular o limite).
\end{itemize}
\end{dicasbox}
\vspace{0.5cm}

\vspace{0.5cm}
\noindent\textbf{Exercício 1.} Seja $f(x) = x^2 + 1$. Calcule:

\begin{enumerate}
    \renewcommand{\theenumi}{\alph{enumi}}
    \renewcommand{\labelenumi}{\theenumi)}
    \item $f'(x)$
    \begin{recursosbox}[title=Recursos Utilizados]
    Utilizaremos a definição de derivada por acréscimo:
    \[
    f'(x) = \lim_{h \to 0} \frac{f(x + h) - f(x)}{h}
    \]
    \textbf{Ferramentas Algébricas:}
    \begin{itemize}[noitemsep, topsep=2pt]
        \item Expansão de binômio: $(x + h)^2 = x^2 + 2hx + h^2$.
        \item Fatoração do quociente incremental para simplificar o fator comum $h$.
    \end{itemize}
    \end{recursosbox}

    Sabemos pela definição formal que:
    \begin{equation}
        f'(x) = \lim_{h \to 0} \frac{f(x + h) - f(x)}{h} \label{eq:def_derivada}
    \end{equation}

    Primeiro, manipulamos o quociente incremental:
    \begin{equation}
        \frac{f(x + h) - f(x)}{h} \label{eq:quociente}
    \end{equation}

    Como foi dado $f(x) = x^2 + 1$, temos:
    \[
    f(x + h) = (x + h)^2 + 1 = x^2 + 2hx + h^2 + 1
    \]
    Substituindo $f(x + h)$ e $f(x)$ em \eqref{eq:quociente} para $h \neq 0$:
    \begin{align*}
        \frac{f(x + h) - f(x)}{h} &= \frac{(x^2 + 2hx + h^2 + 1) - (x^2 + 1)}{h} \\[0.25cm]
        &= \frac{2hx + h^2}{h} \\[0.25cm]
        &= \frac{h(2x + h)}{h} \\[0.25cm]
        &= 2x + h
    \end{align*}

    Daí, aplicando o limite em \eqref{eq:def_derivada} quando $h \to 0$:
    \[
    f'(x) = \lim_{h \to 0} (2x + h) = 2x
    \]
    Portanto:
    \begin{equation*}
        f'(x) = 2x \qedsymbolmath
    \end{equation*}

    \item $f'(0)$
    \setcounter{equation}{0}

    \begin{recursosbox}[title=Recursos Utilizados]
    Utilizaremos a definição formal de derivada no ponto $p = 0$:
    \[
    f'(0) = \lim_{x \to 0} \frac{f(x) - f(0)}{x - 0}
    \]
    \end{recursosbox}

    Pela definição formal de derivada no ponto $p = 0$, temos:
    \begin{equation}
        f'(0) = \lim_{x \to 0} \frac{f(x) - f(0)}{x - 0} \label{eq:def_derivada_ponto0}
    \end{equation}

    Primeiro, manipulamos o quociente incremental para $x \neq 0$:
    \begin{equation}
        \frac{f(x) - f(0)}{x - 0} \label{eq:quociente_ponto0}
    \end{equation}

    Sendo $f(x) = x^2 + 1$, calculamos o valor da função em $p = 0$:
    \[
    f(0) = 0^2 + 1 = 1
    \]

    Substituindo $f(x)$ e $f(0)$ em \eqref{eq:quociente_ponto0}:
    \begin{align*}
        \frac{f(x) - f(0)}{x - 0} &= \frac{(x^2 + 1) - 1}{x} \\[0.25cm]
        &= \frac{x^2}{x} \\[0.25cm]
        &= x
    \end{align*}

    Daí, aplicando o limite em \eqref{eq:def_derivada_ponto0} quando $x \to 0$:
    \[
    f'(0) = \lim_{x \to 0} x = 0
    \]

    Portanto:
    \begin{equation*}
        f'(0) = 0 \qedsymbolmath
    \end{equation*}

    \item $f'(1)$
    \setcounter{equation}{0}

    \begin{recursosbox}[title=Recursos Utilizados]
    Utilizaremos a definição formal de derivada no ponto $p = 1$:
    \[
    f'(1) = \lim_{x \to 1} \frac{f(x) - f(1)}{x - 1}
    \]
    \textbf{Ferramentas Algébricas:}
    \begin{itemize}[noitemsep, topsep=2pt]
        \item Fatoração por diferença de dois quadrados: $x^2 - 1 = (x - 1)(x + 1)$.
    \end{itemize}
    \end{recursosbox}

    Pela definição formal de derivada no ponto $p = 1$, temos:
    \begin{equation}
        f'(1) = \lim_{x \to 1} \frac{f(x) - f(1)}{x - 1} \label{eq:def_derivada_ponto1}
    \end{equation}

    Primeiro, manipulamos o quociente incremental para $x \neq 1$:
    \begin{equation}
        \frac{f(x) - f(1)}{x - 1} \label{eq:quociente_ponto1}
    \end{equation}

    Sendo $f(x) = x^2 + 1$, calculamos o valor da função em $p = 1$:
    \[
    f(1) = 1^2 + 1 = 2
    \]

    Substituindo $f(x)$ e $f(1)$ em \eqref{eq:quociente_ponto1}:
    \begin{align*}
        \frac{f(x) - f(1)}{x - 1} &= \frac{(x^2 + 1) - 2}{x - 1} \\[0.25cm]
        &= \frac{x^2 - 1}{x - 1} \\[0.25cm]
        &= \frac{(x - 1)(x + 1)}{x - 1} \\[0.25cm]
        &= x + 1
    \end{align*}

    Daí, aplicando o limite em \eqref{eq:def_derivada_ponto1} quando $x \to 1$:
    \[
    f'(1) = \lim_{x \to 1} (x + 1) = 1 + 1 = 2
    \]
    Portanto:
    \begin{equation}
        f'(1) = 2 \qedsymbolmath
    \end{equation}
\end{enumerate}

\vspace{0.5cm}
\noindent\textbf{Exercício 2.} Seja $f(x) = 2x$. Pensando geometricamente, qual o valor que você espera para $f'(p)$? Calcule $f'(p)$.
\setcounter{equation}{0}

\begin{recursosbox}[title=Recursos Utilizados]
Utilizaremos a definição formal de derivada no ponto $p$:
\[
f'(p) = \lim_{x \to p} \frac{f(x) - f(p)}{x - p}
\]
e a equação da reta tangente ao gráfico de $f$ em $(p, f(p))$:
\[
y - f(p) = f'(p)(x - p)
\]
\textbf{Ferramentas Algébricas e Geométricas:}
\begin{itemize}[noitemsep, topsep=2pt]
    \item Fatoração por termo comum: $2x - 2p = 2(x - p)$.
    \item Interpretação geométrica: $f(x) = 2x$ representa uma reta com inclinação constante $m = 2$.
\end{itemize}
\end{recursosbox}

\textbf{Passo 1: Cálculo de $f'(p)$ pela definição formal.}

Pela definição formal de derivada no ponto $p$:
\begin{equation}
    f'(p) = \lim_{x \to p} \frac{f(x) - f(p)}{x - p} \label{eq:def_derivada_ex2}
\end{equation}

Primeiro, manipulamos o quociente incremental para $x \neq p$:
\begin{equation}
    \frac{f(x) - f(p)}{x - p} \label{eq:quociente_ex2}
\end{equation}

Sendo $f(x) = 2x$, temos $f(p) = 2p$. Substituindo em \eqref{eq:quociente_ex2}:
\begin{align*}
    \frac{f(x) - f(p)}{x - p} &= \frac{2x - 2p}{x - p} \\[0.25cm]
    &= \frac{2(x - p)}{x - p} \\[0.25cm]
    &= 2
\end{align*}

Aplicando o limite em \eqref{eq:def_derivada_ex2} quando $x \to p$:
\[
f'(p) = \lim_{x \to p} 2 = 2
\]

Portanto:
\begin{equation*}
    f'(p) = 2 \qedsymbolmath
\end{equation*}

\vspace{0.3cm}
\textbf{Passo 2: Discussão geométrica do resultado.}

Geometricamente, a derivada $f'(p)$ representa a taxa de variação instantânea de $f$ em relação a $x$, isto é, a inclinação (coeficiente angular) da reta tangente à curva no ponto $(p, f(p))$.

Como $f(x) = 2x$ já é uma função linear cujo gráfico é uma reta de equação $y = 2x$, a sua inclinação é constante em todos os pontos e igual a $2$. Consequentemente:
\begin{itemize}[noitemsep, topsep=2pt]
    \item Espera-se geometricamente que $f'(p) = 2$ para qualquer $p \in \mathbb{R}$.
\end{itemize}

\vspace{0.3cm}
\textbf{Passo 3: Determinação da equação da reta tangente em $(p, f(p))$.}

A equação da reta tangente ao gráfico de $f$ no ponto $(p, f(p))$ é dada por:
\begin{equation}
    y - f(p) = f'(p)(x - p) \label{eq:reta_tangente_ex2}
\end{equation}

Sabendo que $f(p) = 2p$ e $f'(p) = 2$, substituímos esses valores em \eqref{eq:reta_tangente_ex2}:
\begin{align*}
    y - 2p &= 2(x - p) \\[0.25cm]
    y - 2p &= 2x - 2p \\[0.25cm]
    y &= 2x
\end{align*}

\vspace{0.3cm}
\textbf{Passo 4: Discussão gráfica da reta tangente.}

A reta tangente obtida tem por equação $y = 2x$, que coincide exatamente com a própria função $f(x) = 2x$.

Graficamente, isso expressa uma propriedade fundamental das funções afins: a melhor aproximação linear local para uma reta em qualquer um de seus pontos é a própria reta. Assim, em qualquer ponto $(p, 2p)$, a reta tangente é coincidente com o gráfico da função em toda a sua extensão.

\begin{center}
\begin{tikzpicture}[>=latex, scale=0.95]
    % Eixos coordenados
    \draw[->, thick] (-2.2, 0) -- (3.2, 0) node[below] {$x$};
    \draw[->, thick] (0, -2.5) -- (0, 4.2) node[left] {$y$};
    \node[below left=2pt] at (0, 0) {$0$};

    % Reta tangente coincidente
    \draw[line width=3pt, blue!20] (-1.2, -2.4) -- (2.0, 4.0);

    % Gráfico de f(x) = 2x
    \draw[line width=1.2pt, blue!70!black] (-1.2, -2.4) -- (2.0, 4.0);

    % Ponto genérico (p, 2p)
    \coordinate (P) at (1.2, 2.4);

    % Linhas tracejadas de projeção
    \draw[dashed, gray!80] (1.2, 0) node[below, black] {$p$} -- (P);
    \draw[dashed, gray!80] (0, 2.4) node[left, black] {$f(p) = 2p$} -- (P);

    % Marcação do ponto
    \fill[blue!70!black] (P) circle (3pt);
    \node[above left=3pt] at (P) {$(p, 2p)$};

    % Rótulos explicativos
    \node[right, blue!70!black] at (1.8, 3.6) {$f(x) = 2x$};
    \node[right, blue!50!black] at (1.8, 3.1) {\footnotesize Reta Tangente: $y = 2x$};
\end{tikzpicture}
\end{center}

\vspace{0.5cm}
\noindent\textbf{Exercício 3.} Seja $f(x) = 3x + 2$. Calcule:
\begin{enumerate}
    \renewcommand{\theenumi}{\alph{enumi}}
    \renewcommand{\labelenumi}{\theenumi)}
    \item $f'(2)$
    \item $f'(0)$
    \item $f'(x)$
\end{enumerate}
\setcounter{equation}{0}

\begin{recursosbox}[title=Estratégia e Recursos Utilizados]
\textbf{Estratégia Geral:} Em vez de calcular o limite da definição separadamente para cada ponto numérico, determinaremos primeiramente a derivada $f'(p)$ em um ponto genérico $p \in \operatorname{D}_f$ (ou a função derivada $f'(x)$). Em seguida, utilizaremos essa expressão geral para avaliar diretamente os itens pontuais por substituição simples.

\vspace{0.2cm}
\textbf{Definição Formal de Derivada em $p$:}
\[
f'(p) = \lim_{x \to p} \frac{f(x) - f(p)}{x - p}
\]
\textbf{Ferramentas Algébricas:}
\begin{itemize}[noitemsep, topsep=2pt]
    \item Fatoração por termo comum: $3x - 3p = 3(x - p)$.
    \item Cancelamento do fator indeterminante $(x - p)$ para $x \neq p$.
\end{itemize}
\end{recursosbox}

\textbf{Determinação da Derivada Geral em um Ponto $p \in \mathbb{R}$:}

Pela definição formal de derivada num ponto arbitrário $p$:
\begin{equation}
    f'(p) = \lim_{x \to p} \frac{f(x) - f(p)}{x - p} \label{eq:def_ex3}
\end{equation}

Para $x \neq p$, manipulamos algebricamente o quociente incremental:
\begin{equation}
    \frac{f(x) - f(p)}{x - p} \label{eq:quociente_ex3}
\end{equation}

Como $f(x) = 3x + 2$, temos $f(p) = 3p + 2$. Substituindo em \eqref{eq:quociente_ex3}:
\begin{align*}
    \frac{f(x) - f(p)}{x - p} &= \frac{(3x + 2) - (3p + 2)}{x - p} \\[0.25cm]
    &= \frac{3x - 3p}{x - p} \\[0.25cm]
    &= \frac{3(x - p)}{x - p} \\[0.25cm]
    &= 3
\end{align*}

Aplicando o limite em \eqref{eq:def_ex3} quando $x \to p$:
\[
f'(p) = \lim_{x \to p} 3 = 3
\]
Como $p$ é um elemento arbitrário do domínio $\mathbb{R}$, concluímos que a derivada é constante e independe do ponto escolhido:
\begin{equation}
    f'(p) = 3 \quad \forall p \in \mathbb{R} \label{eq:resultado_geral_ex3}
\end{equation}

\vspace{0.3cm}
\textbf{Resolução dos Itens por Substituição Direta:}

\begin{enumerate}
    \renewcommand{\theenumi}{\alph{enumi}}
    \renewcommand{\labelenumi}{\theenumi)}
    \item $f'(2)$
    
    Fazendo $p = 2$ no resultado geral \eqref{eq:resultado_geral_ex3}:
    \begin{equation*}
        f'(2) = 3 \qedsymbolmath
    \end{equation*}

    \item $f'(0)$
    
    Fazendo $p = 0$ no resultado geral \eqref{eq:resultado_geral_ex3}:
    \begin{equation*}
        f'(0) = 3 \qedsymbolmath
    \end{equation*}

    \item $f'(x)$
    
    Sendo a variável independente indicada por $x$, a função derivada associa a cada ponto $x \in \mathbb{R}$ o valor de sua derivada obtido em \eqref{eq:resultado_geral_ex3}:
    \begin{equation*}
        f'(x) = 3 \qedsymbolmath
    \end{equation*}
\end{enumerate}

\vspace{0.5cm}
\noindent\textbf{Exercício 4.} Calcule $f'(p)$, pela definição, sendo dados:
\begin{enumerate}[label=\alph*)]
    \item $f(x) = x^2 + x$ e $p = 1$
    \item $f(x) = \sqrt{x}$ e $p = 4$
    \item $f(x) = 5x - 3$ e $p = -3$
    \item $f(x) = \dfrac{1}{x}$ e $p = 1$
    \item $f(x) = \sqrt{x}$ e $p = 3$
    \item $f(x) = \dfrac{1}{x^2}$ e $p = 2$
    \item $f(x) = 2x^3 - x^2$ e $p = 1$
    \item $f(x) = \sqrt[3]{x}$ e $p = 2$
\end{enumerate}
\setcounter{equation}{0}

\begin{recursosbox}[title=Estratégia Geral e Recursos Algébricos]
\textbf{Definição Formal de Derivada no Ponto $p$:}
\[
f'(p) = \lim_{x \to p} \frac{f(x) - f(p)}{x - p}
\]
Como a avaliação direta em $x = p$ sempre conduz à forma indeterminada $\left[\frac{0}{0}\right]$, o objetivo central é manipular algebricamente o quociente incremental $\dfrac{f(x) - f(p)}{x - p}$ de modo a simplificar o fator $(x - p)$ para $x \neq p$:
\begin{itemize}[noitemsep, topsep=2pt]
    \item \textbf{Polinômios:} Fatoração direta por agrupamento ou raízes para evidenciar o fator $(x - p)$.
    \item \textbf{Funções Racionais:} Operação com frações (denominador comum) e simplificação direta de termos semelhantes.
    \item \textbf{Raízes Quadradas:} Multiplicação pelo conjugado utilizando a diferença de quadrados:
    \[
    (\sqrt{A} - \sqrt{B})(\sqrt{A} + \sqrt{B}) = A - B
    \]
    \item \textbf{Raízes Cúbicas:} Utilização do produto notável da diferença de cubos:
    \[
    A - B = (\sqrt[3]{A} - \sqrt[3]{B})\left(\sqrt[3]{A^2} + \sqrt[3]{AB} + \sqrt[3]{B^2}\right)
    \]
    \item \textbf{Racionalização do Resultado Final (Radicais no Denominador):} Para simplificar respostas numéricas finais com frações irredutíveis:
    \begin{itemize}[noitemsep, topsep=1pt]
        \item \textit{Raiz quadrada:} Multiplica-se numerador e denominador por $\sqrt{a}$:
        \[
        \frac{1}{k\sqrt{a}} = \frac{\sqrt{a}}{k \cdot a} \quad \left(\text{ex.: } \frac{1}{2\sqrt{3}} = \frac{\sqrt{3}}{6}\right)
        \]
        \item \textit{Raiz de ordem $n$ ($\sqrt[n]{a^m}$ com $m < n$):} Multiplica-se por $\sqrt[n]{a^{n - m}}$ para completar a potência inteira $a^n$:
        \[
        \frac{1}{k\sqrt[n]{a^m}} = \frac{\sqrt[n]{a^{n - m}}}{k \cdot a} \quad \left(\text{ex.: } \frac{1}{3\sqrt[3]{4}} = \frac{1}{3\sqrt[3]{2^2}} = \frac{\sqrt[3]{2}}{3 \cdot 2} = \frac{\sqrt[3]{2}}{6}\right)
        \]
    \end{itemize}
\end{itemize}
\end{recursosbox}

\vspace{0.4cm}
\textbf{Resoluções Passo a Passo:}

\begin{enumerate}[label=\alph*)]
    % Item a
    \item $f(x) = x^2 + x$ e $p = 1$

    Calculando o valor de $f$ em $p = 1$:
    \[
    f(1) = 1^2 + 1 = 2
    \]
    Para $x \neq 1$, simplificamos o quociente incremental:
    \begin{align*}
        \frac{f(x) - f(1)}{x - 1} &= \frac{x^2 + x - 2}{x - 1} \\[0.25cm]
        &= \frac{(x - 1)(x + 2)}{x - 1} \\[0.25cm]
        &= x + 2
    \end{align*}
    Calculando o limite pela definição:
    \[
    f'(1) = \lim_{x \to 1} (x + 2) = 1 + 2 = 3
    \]
    \begin{equation*}
        f'(1) = 3 \qedsymbolmath
    \end{equation*}

    % Item b
    \item $f(x) = \sqrt{x}$ e $p = 4$

    Calculando o valor de $f$ em $p = 4$:
    \[
    f(4) = \sqrt{4} = 2
    \]
    Para $x \neq 4$ com $x \ge 0$, multiplicamos numerador e denominador pelo conjugado $(\sqrt{x} + 2)$:
    \begin{align*}
        \frac{f(x) - f(4)}{x - 4} &= \frac{\sqrt{x} - 2}{x - 4} \\[0.25cm]
        &= \frac{(\sqrt{x} - 2)(\sqrt{x} + 2)}{(x - 4)(\sqrt{x} + 2)} \\[0.25cm]
        &= \frac{x - 4}{(x - 4)(\sqrt{x} + 2)} \\[0.25cm]
        &= \frac{1}{\sqrt{x} + 2}
    \end{align*}
    Passando ao limite quando $x \to 4$:
    \[
    f'(4) = \lim_{x \to 4} \frac{1}{\sqrt{x} + 2} = \frac{1}{\sqrt{4} + 2} = \frac{1}{4}
    \]
    \begin{equation*}
        f'(4) = \frac{1}{4} \qedsymbolmath
    \end{equation*}

    % Item c
    \item $f(x) = 5x - 3$ e $p = -3$

    Calculando o valor de $f$ em $p = -3$:
    \[
    f(-3) = 5(-3) - 3 = -18
    \]
    Para $x \neq -3$:
    \begin{align*}
        \frac{f(x) - f(-3)}{x - (-3)} &= \frac{(5x - 3) - (-18)}{x + 3} \\[0.25cm]
        &= \frac{5x + 15}{x + 3} \\[0.25cm]
        &= \frac{5(x + 3)}{x + 3} = 5
    \end{align*}
    Passando ao limite quando $x \to -3$:
    \[
    f'(-3) = \lim_{x \to -3} 5 = 5
    \]
    \begin{equation*}
        f'(-3) = 5 \qedsymbolmath
    \end{equation*}

    % Item d
    \item $f(x) = \dfrac{1}{x}$ e $p = 1$

    Calculando o valor de $f$ em $p = 1$:
    \[
    f(1) = \frac{1}{1} = 1
    \]
    Para $x \neq 1$ e $x \neq 0$:
    \begin{align*}
        \frac{f(x) - f(1)}{x - 1} &= \frac{\dfrac{1}{x} - 1}{x - 1} \\[0.25cm]
        &= \frac{\dfrac{1 - x}{x}}{x - 1} \\[0.25cm]
        &= \frac{-(x - 1)}{x(x - 1)} \\[0.25cm]
        &= -\frac{1}{x}
    \end{align*}
    Passando ao limite quando $x \to 1$:
    \[
    f'(1) = \lim_{x \to 1} \left(-\frac{1}{x}\right) = -1
    \]
    \begin{equation*}
        f'(1) = -1 \qedsymbolmath
    \end{equation*}

    % Item e
    \item $f(x) = \sqrt{x}$ e $p = 3$

    Calculando o valor de $f$ em $p = 3$:
    \[
    f(3) = \sqrt{3}
    \]
    Para $x \neq 3$ e $x \ge 0$, multiplicamos pelo conjugado $(\sqrt{x} + \sqrt{3})$:
    \begin{align*}
        \frac{f(x) - f(3)}{x - 3} &= \frac{\sqrt{x} - \sqrt{3}}{x - 3} \\[0.25cm]
        &= \frac{(\sqrt{x} - \sqrt{3})(\sqrt{x} + \sqrt{3})}{(x - 3)(\sqrt{x} + \sqrt{3})} \\[0.25cm]
        &= \frac{x - 3}{(x - 3)(\sqrt{x} + \sqrt{3})} \\[0.25cm]
        &= \frac{1}{\sqrt{x} + \sqrt{3}}
    \end{align*}
    Passando ao limite quando $x \to 3$:
    \[
    f'(3) = \lim_{x \to 3} \frac{1}{\sqrt{x} + \sqrt{3}} = \frac{1}{\sqrt{3} + \sqrt{3}} = \frac{1}{2\sqrt{3}} = \frac{\sqrt{3}}{6}
    \]
    \begin{equation*}
        f'(3) = \frac{\sqrt{3}}{6} \qedsymbolmath
    \end{equation*}

    % Item f
    \item $f(x) = \dfrac{1}{x^2}$ e $p = 2$

    Calculando o valor de $f$ em $p = 2$:
    \[
    f(2) = \frac{1}{2^2} = \frac{1}{4}
    \]
    Para $x \neq 2$ e $x \neq 0$:
    \begin{align*}
        \frac{f(x) - f(2)}{x - 2} &= \frac{\dfrac{1}{x^2} - \dfrac{1}{4}}{x - 2} \\[0.25cm]
        &= \frac{\dfrac{4 - x^2}{4x^2}}{x - 2} \\[0.25cm]
        &= \frac{-(x - 2)(x + 2)}{4x^2(x - 2)} \\[0.25cm]
        &= -\frac{x + 2}{4x^2}
    \end{align*}
    Passando ao limite quando $x \to 2$:
    \[
    f'(2) = \lim_{x \to 2} \left(-\frac{x + 2}{4x^2}\right) = -\frac{2 + 2}{4(2^2)} = -\frac{4}{16} = -\frac{1}{4}
    \]
    \begin{equation*}
        f'(2) = -\frac{1}{4} \qedsymbolmath
    \end{equation*}

    % Item g
    \item $f(x) = 2x^3 - x^2$ e $p = 1$

    Calculando o valor de $f$ em $p = 1$:
    \[
    f(1) = 2(1)^3 - 1^2 = 1
    \]
    Para $x \neq 1$, manipulamos o quociente incremental:
    \[
    \frac{f(x) - f(1)}{x - 1} = \frac{2x^3 - x^2 - 1}{x - 1}
    \]
    Fatorando o polinômio do numerador $2x^3 - x^2 - 1$ por $(x - 1)$:
    \[
    2x^3 - x^2 - 1 = (x - 1)(2x^2 + x + 1)
    \]
    Portanto:
    \begin{align*}
        \frac{f(x) - f(1)}{x - 1} &= \frac{(x - 1)(2x^2 + x + 1)}{x - 1} = 2x^2 + x + 1
    \end{align*}
    Passando ao limite quando $x \to 1$:
    \[
    f'(1) = \lim_{x \to 1} (2x^2 + x + 1) = 2(1)^2 + 1 + 1 = 4
    \]
    \begin{equation*}
        f'(1) = 4 \qedsymbolmath
    \end{equation*}

    % Item h
    \item $f(x) = \sqrt[3]{x}$ e $p = 2$

    Calculando o valor de $f$ em $p = 2$:
    \[
    f(2) = \sqrt[3]{2}
    \]
    Para $x \neq 2$:
    \[
    \frac{f(x) - f(2)}{x - 2} = \frac{\sqrt[3]{x} - \sqrt[3]{2}}{x - 2}
    \]
    Utilizando a identidade algébrica da diferença de cubos $A - B = (\sqrt[3]{A} - \sqrt[3]{B})\left(\sqrt[3]{A^2} + \sqrt[3]{AB} + \sqrt[3]{B^2}\right)$ no denominador para simplificar o fator $(\sqrt[3]{x} - \sqrt[3]{2})$:
    \begin{align*}
        \frac{f(x) - f(2)}{x - 2} &= \frac{\sqrt[3]{x} - \sqrt[3]{2}}{(\sqrt[3]{x} - \sqrt[3]{2})\left(\sqrt[3]{x^2} + \sqrt[3]{2x} + \sqrt[3]{4}\right)} \\[0.25cm]
        &= \frac{1}{\sqrt[3]{x^2} + \sqrt[3]{2x} + \sqrt[3]{4}}
    \end{align*}
    Passando ao limite quando $x \to 2$:
    \[
    f'(2) = \lim_{x \to 2} \frac{1}{\sqrt[3]{x^2} + \sqrt[3]{2x} + \sqrt[3]{4}} = \frac{1}{\sqrt[3]{4} + \sqrt[3]{4} + \sqrt[3]{4}} = \frac{1}{3\sqrt[3]{4}} = \frac{\sqrt[3]{2}}{6}
    \]
    \begin{equation*}
        f'(2) = \frac{\sqrt[3]{2}}{6} \qedsymbolmath
    \end{equation*}
\end{enumerate}

\vspace{0.5cm}
\noindent\textbf{Exercício 5.} Determine a equação da reta tangente em $(p, f(p))$ sendo dados:
\begin{enumerate}
    \renewcommand{\theenumi}{\alph{enumi}}
    \renewcommand{\labelenumi}{\theenumi)}
    \item $f(x) = x^2$ e $p = 2$
    \item $f(x) = \dfrac{1}{x}$ e $p = 2$
    \item $f(x) = \sqrt{x}$ e $p = 9$
    \item $f(x) = x^2 - x$ e $p = 1$
\end{enumerate}
\setcounter{equation}{0}

\begin{recursosbox}[title=Estratégia Geral e Recursos Algébricos]
\textbf{1. Definição Formal de Derivada no Ponto $p$:}
\[
f'(p) = \lim_{x \to p} \frac{f(x) - f(p)}{x - p}
\]
\textbf{2. Equação da Reta Tangente:}
O valor $f'(p)$ corresponde geometricamente ao coeficiente angular ($m$) da reta tangente à curva $y = f(x)$ no ponto de tangência $(p, f(p))$. A equação da reta tangente é dada por:
\[
y - f(p) = f'(p)(x - p)
\]
\textbf{3. Roteiro de Resolução para Cada Item:}
\begin{itemize}[noitemsep, topsep=2pt]
    \item \textbf{Passo 1:} Determinar o ponto de tangência calculando $f(p)$.
    \item \textbf{Passo 2:} Calcular o coeficiente angular $f'(p)$ pela definição de limite.
    \item \textbf{Passo 3:} Substituir $(p, f(p))$ e $f'(p)$ na equação da reta e simplificar para a forma reduzida $y = mx + b$.
\end{itemize}
\end{recursosbox}

\vspace{0.4cm}
\textbf{Resolução Passo a Passo:}

\begin{enumerate}
    \renewcommand{\theenumi}{\alph{enumi}}
    \renewcommand{\labelenumi}{\theenumi)}
    
    % Item a
    \item $f(x) = x^2$ e $p = 2$
    
    \textbf{Passo 1: Determinar o ponto de tangência.}
    \[
    f(2) = 2^2 = 4 \implies (p, f(p)) = (2, 4)
    \]

    \textbf{Passo 2: Calcular o coeficiente angular $f'(2)$ pela definição.}
    Para $x \neq 2$:
    \begin{align*}
        \frac{f(x) - f(2)}{x - 2} &= \frac{x^2 - 4}{x - 2} \\[0.25cm]
        &= \frac{(x - 2)(x + 2)}{x - 2} \\[0.25cm]
        &= x + 2
    \end{align*}
    Passando ao limite quando $x \to 2$:
    \[
    f'(2) = \lim_{x \to 2} (x + 2) = 2 + 2 = 4
    \]

    \textbf{Passo 3: Equação da reta tangente.}

    Substituindo $(p, f(p)) = (2, 4)$ e $f'(2) = 4$:
    \begin{align*}
        y - 4 &= 4(x - 2) \\[0.25cm]
        y - 4 &= 4x - 8 \\[0.25cm]
        y &= 4x - 4
    \end{align*}

    Portanto:
    \begin{equation*}
        y = 4x - 4 \qedsymbolmath
    \end{equation*}

    % Item b
    \item $f(x) = \dfrac{1}{x}$ e $p = 2$
    
    \textbf{Passo 1: Determinar o ponto de tangência.}
    \[
    f(2) = \frac{1}{2} \implies (p, f(p)) = \left(2, \frac{1}{2}\right)
    \]

    \textbf{Passo 2: Calcular o coeficiente angular $f'(2)$ pela definição.}

    Para $x \neq 2$ e $x \neq 0$:
    \begin{align*}
        \frac{f(x) - f(2)}{x - 2} &= \frac{\dfrac{1}{x} - \dfrac{1}{2}}{x - 2} \\[0.25cm]
        &= \frac{\dfrac{2 - x}{2x}}{x - 2} \\[0.25cm]
        &= \frac{-(x - 2)}{2x(x - 2)} \\[0.25cm]
        &= -\frac{1}{2x}
    \end{align*}
    Passando ao limite quando $x \to 2$:
    \[
    f'(2) = \lim_{x \to 2} \left(-\frac{1}{2x}\right) = -\frac{1}{2(2)} = -\frac{1}{4}
    \]

    \textbf{Passo 3: Equação da reta tangente.}

    Substituindo $(p, f(p)) = \left(2, \dfrac{1}{2}\right)$ e $f'(2) = -\dfrac{1}{4}$:
    \begin{align*}
        y - \frac{1}{2} &= -\frac{1}{4}(x - 2) \\[0.25cm]
        y - \frac{1}{2} &= -\frac{1}{4}x + \frac{1}{2} \\[0.25cm]
        y &= -\frac{1}{4}x + 1
    \end{align*}

    Portanto:
    \begin{equation*}
        y = -\frac{1}{4}x + 1 \qedsymbolmath
    \end{equation*}

    % Item c
    \item $f(x) = \sqrt{x}$ e $p = 9$
    
    \textbf{Passo 1: Determinar o ponto de tangência.}
    \[
    f(9) = \sqrt{9} = 3 \implies (p, f(p)) = (9, 3)
    \]

    \textbf{Passo 2: Calcular o coeficiente angular $f'(9)$ pela definição.}

    Para $x \neq 9$ com $x \ge 0$, multiplicamos pelo conjugado $(\sqrt{x} + 3)$:
    \begin{align*}
        \frac{f(x) - f(9)}{x - 9} &= \frac{\sqrt{x} - 3}{x - 9} \\[0.25cm]
        &= \frac{(\sqrt{x} - 3)(\sqrt{x} + 3)}{(x - 9)(\sqrt{x} + 3)} \\[0.25cm]
        &= \frac{x - 9}{(x - 9)(\sqrt{x} + 3)} \\[0.25cm]
        &= \frac{1}{\sqrt{x} + 3}
    \end{align*}
    Passando ao limite quando $x \to 9$:
    \[
    f'(9) = \lim_{x \to 9} \frac{1}{\sqrt{x} + 3} = \frac{1}{\sqrt{9} + 3} = \frac{1}{6}
    \]

    \textbf{Passo 3: Equação da reta tangente.}

    Substituindo $(p, f(p)) = (9, 3)$ e $f'(9) = \dfrac{1}{6}$:
    \begin{align*}
        y - 3 &= \frac{1}{6}(x - 9) \\[0.25cm]
        y - 3 &= \frac{1}{6}x - \frac{3}{2} \\[0.25cm]
        y &= \frac{1}{6}x + \frac{3}{2}
    \end{align*}

    Portanto:
    \begin{equation*}
        y = \frac{1}{6}x + \frac{3}{2} \qedsymbolmath
    \end{equation*}

    % Item d
    \item $f(x) = x^2 - x$ e $p = 1$
    
    \textbf{Passo 1: Determinar o ponto de tangência.}
    \[
    f(1) = 1^2 - 1 = 0 \implies (p, f(p)) = (1, 0)
    \]

    \textbf{Passo 2: Calcular o coeficiente angular $f'(1)$ pela definição.}

    Para $x \neq 1$:
    \begin{align*}
        \frac{f(x) - f(1)}{x - 1} &= \frac{(x^2 - x) - 0}{x - 1} \\[0.25cm]
        &= \frac{x(x - 1)}{x - 1} \\[0.25cm]
        &= x
    \end{align*}
    Passando ao limite quando $x \to 1$:
    \[
    f'(1) = \lim_{x \to 1} x = 1
    \]

    \textbf{Passo 3: Equação da reta tangente.}

    Substituindo $(p, f(p)) = (1, 0)$ e $f'(1) = 1$:
    \begin{align*}
        y - 0 &= 1(x - 1) \\[0.25cm]
        y &= x - 1
    \end{align*}

    Portanto:
    \begin{equation*}
        y = x - 1 \qedsymbolmath
    \end{equation*}
\end{enumerate}

\vspace{0.5cm}
\noindent\textbf{Exercício 6.}
\begin{enumerate}
    \renewcommand{\theenumi}{\alph{enumi}}
    \renewcommand{\labelenumi}{\theenumi)}
    \item 
    \item 
    \item 
    \item 
    \item 
    \item 
    \item 
    \item 
\end{enumerate}

\vspace{0.5cm}
\noindent\textbf{Exercício 7.}

\vspace{0.5cm}
\noindent\textbf{Exercício 8.}

\vspace{0.5cm}
\noindent\textbf{Exercício 9.}

\vspace{0.5cm}
\noindent\textbf{Exercício 10.}

\vspace{0.5cm}
\noindent\textbf{Exercício 11.}

\vspace{0.5cm}
\noindent\textbf{Exercício 12.}

\vspace{0.5cm}
\noindent\textbf{Exercício 13.}

\vspace{0.5cm}
\noindent\textbf{Exercício 14.}

\vspace{0.5cm}
\noindent\textbf{Exercício 15.}
\begin{enumerate}
    \renewcommand{\theenumi}{\alph{enumi}}
    \renewcommand{\labelenumi}{\theenumi)}
    \item 
    \item 
\end{enumerate}

\vspace{0.5cm}
\noindent\textbf{Exercício 16.}
\begin{enumerate}
    \renewcommand{\theenumi}{\alph{enumi}}
    \renewcommand{\labelenumi}{\theenumi)}
    \item 
    \item 
\end{enumerate}

\vspace{0.5cm}
\noindent\textbf{Exercício 17.}
\begin{enumerate}
    \renewcommand{\theenumi}{\alph{enumi}}
    \renewcommand{\labelenumi}{\theenumi)}
    \item 
    \item 
\end{enumerate}

\vspace{0.5cm}
\noindent\textbf{Exercício 18.}

\vspace{0.5cm}
\noindent\textbf{Exercício 19.}

\subsection{Derivadas de $x^n$ e $\sqrt[n]{x}$}

\vspace{0.5cm}
\noindent\textbf{Exercício 1.}
\begin{enumerate}
    \renewcommand{\theenumi}{\alph{enumi}}
    \renewcommand{\labelenumi}{\theenumi)}
    \item 
    \item 
    \item 
\end{enumerate}

\vspace{0.5cm}
\noindent\textbf{Exercício 2.}
\begin{enumerate}
    \renewcommand{\theenumi}{\alph{enumi}}
    \renewcommand{\labelenumi}{\theenumi)}
    \item 
    \item 
    \item 
    \item 
    \item 
    \item 
    \item 
    \item 
\end{enumerate}

\vspace{0.5cm}
\noindent\textbf{Exercício 3.}

\vspace{0.5cm}
\noindent\textbf{Exercício 4.}

\vspace{0.5cm}
\noindent\textbf{Exercício 5.}
\begin{enumerate}
    \renewcommand{\theenumi}{\alph{enumi}}
    \renewcommand{\labelenumi}{\theenumi)}
    \item 
    \item 
    \item 
\end{enumerate}

\vspace{0.5cm}
\noindent\textbf{Exercício 6.}
\begin{enumerate}
    \renewcommand{\theenumi}{\alph{enumi}}
    \renewcommand{\labelenumi}{\theenumi)}
    \item 
    \item 
    \item 
    \item 
\end{enumerate}

\vspace{0.5cm}
\noindent\textbf{Exercício 7.}

\vspace{0.5cm}
\noindent\textbf{Exercício 8.}

\vspace{0.5cm}
\noindent\textbf{Exercício 9.}

\subsection{Derivadas de $e^x$ e $\ln(x)$}

\vspace{0.5cm}
\noindent\textbf{Exercício 1.}

\vspace{0.5cm}
\noindent\textbf{Exercício 2.}

\vspace{0.5cm}
\noindent\textbf{Exercício 3.}

\vspace{0.5cm}
\noindent\textbf{Exercício 4.}
\begin{enumerate}
    \renewcommand{\theenumi}{\alph{enumi}}
    \renewcommand{\labelenumi}{\theenumi)}
    \item 
    \item 
    \item 
    \item 
\end{enumerate}

\vspace{0.5cm}
\noindent\textbf{Exercício 5.}

\vspace{0.5cm}
\noindent\textbf{Exercício 6.}
\begin{enumerate}
    \renewcommand{\theenumi}{\alph{enumi}}
    \renewcommand{\labelenumi}{\theenumi)}
    \item 
    \item 
    \item 
    \item 
\end{enumerate}

\subsection{Derivadas de Funções Trigonométricas}

\vspace{0.5cm}
\noindent\textbf{Exercício 1.}
\begin{enumerate}
    \renewcommand{\theenumi}{\alph{enumi}}
    \renewcommand{\labelenumi}{\theenumi)}
    \item 
    \item 
\end{enumerate}

\vspace{0.5cm}
\noindent\textbf{Exercício 2.}

\vspace{0.5cm}
\noindent\textbf{Exercício 3.}
\begin{enumerate}
    \renewcommand{\theenumi}{\alph{enumi}}
    \renewcommand{\labelenumi}{\theenumi)}
    \item 
    \item 
    \item 
    \item 
\end{enumerate}

\vspace{0.5cm}
\noindent\textbf{Exercício 4.}
\begin{enumerate}
    \renewcommand{\theenumi}{\alph{enumi}}
    \renewcommand{\labelenumi}{\theenumi)}
    \item 
    \item 
\end{enumerate}

\vspace{0.5cm}
\noindent\textbf{Exercício 5.}

\vspace{0.5cm}
\noindent\textbf{Exercício 6.}
\begin{enumerate}
    \renewcommand{\theenumi}{\alph{enumi}}
    \renewcommand{\labelenumi}{\theenumi)}
    \item 
    \item 
\end{enumerate}

\vspace{0.5cm}
\noindent\textbf{Exercício 7.}
\begin{enumerate}
    \renewcommand{\theenumi}{\alph{enumi}}
    \renewcommand{\labelenumi}{\theenumi)}
    \item 
    \item 
\end{enumerate}

\subsection{Derivabilidade e Continuidade}

\vspace{0.5cm}
\noindent\textbf{Exercício 1.}
\begin{enumerate}
    \renewcommand{\theenumi}{\alph{enumi}}
    \renewcommand{\labelenumi}{\theenumi)}
    \item 
    \item 
\end{enumerate}

\vspace{0.5cm}
\noindent\textbf{Exercício 2.}
\begin{enumerate}
    \renewcommand{\theenumi}{\alph{enumi}}
    \renewcommand{\labelenumi}{\theenumi)}
    \item 
    \item 
\end{enumerate}

\vspace{0.5cm}
\noindent\textbf{Exercício 3.}
\begin{enumerate}
    \renewcommand{\theenumi}{\alph{enumi}}
    \renewcommand{\labelenumi}{\theenumi)}
    \item 
    \item 
\end{enumerate}

\subsection{Regras de Derivação}

\vspace{0.5cm}
\noindent\textbf{Exercício 1.}
\begin{enumerate}
    \renewcommand{\theenumi}{\alph{enumi}}
    \renewcommand{\labelenumi}{\theenumi)}
    \item 
    \item 
    \item 
    \item 
    \item 
    \item 
    \item 
    \item 
    \item 
    \item 
    \item 
    \item 
    \item 
\end{enumerate}

\vspace{0.5cm}
\noindent\textbf{Exercício 2.}

\vspace{0.5cm}
\noindent\textbf{Exercício 3.}
\begin{enumerate}
    \renewcommand{\theenumi}{\alph{enumi}}
    \renewcommand{\labelenumi}{\theenumi)}
    \item 
    \item 
\end{enumerate}

\vspace{0.5cm}
\noindent\textbf{Exercício 4.}
\begin{enumerate}
    \renewcommand{\theenumi}{\alph{enumi}}
    \renewcommand{\labelenumi}{\theenumi)}
    \item 
    \item 
    \item 
\end{enumerate}

\vspace{0.5cm}
\noindent\textbf{Exercício 5.}

\vspace{0.5cm}
\noindent\textbf{Exercício 6.}
\begin{enumerate}
    \renewcommand{\theenumi}{\alph{enumi}}
    \renewcommand{\labelenumi}{\theenumi)}
    \item 
    \item 
    \item 
\end{enumerate}

\vspace{0.5cm}
\noindent\textbf{Exercício 7.}
\begin{enumerate}
    \renewcommand{\theenumi}{\alph{enumi}}
    \renewcommand{\labelenumi}{\theenumi)}
    \item 
    \item 
    \item 
    \item 
    \item 
    \item 
    \item 
    \item 
\end{enumerate}

\vspace{0.5cm}
\noindent\textbf{Exercício 8.}
\begin{enumerate}
    \renewcommand{\theenumi}{\alph{enumi}}
    \renewcommand{\labelenumi}{\theenumi)}
    \item 
    \item 
    \item 
    \item 
\end{enumerate}

\vspace{0.5cm}
\noindent\textbf{Exercício 9.}
\begin{enumerate}
    \renewcommand{\theenumi}{\alph{enumi}}
    \renewcommand{\labelenumi}{\theenumi)}
    \item 
    \item 
    \item 
    \item 
    \item 
    \item 
    \item 
    \item 
    \item 
    \item 
    \item 
    \item 
    \item 
    \item 
    \item 
    \item 
    \item 
    \item 
\end{enumerate}

\vspace{0.5cm}
\noindent\textbf{Exercício 10.}
\begin{enumerate}
    \renewcommand{\theenumi}{\alph{enumi}}
    \renewcommand{\labelenumi}{\theenumi)}
    \item 
    \item 
    \item 
    \item 
\end{enumerate}

\vspace{0.5cm}
\noindent\textbf{Exercício 11.}
\begin{enumerate}
    \renewcommand{\theenumi}{\alph{enumi}}
    \renewcommand{\labelenumi}{\theenumi)}
    \item 
    \item 
\end{enumerate}

\vspace{0.5cm}
\noindent\textbf{Exercício 12.}
\begin{enumerate}
    \renewcommand{\theenumi}{\alph{enumi}}
    \renewcommand{\labelenumi}{\theenumi)}
    \item 
    \item 
    \item 
    \item 
    \item 
    \item 
    \item 
    \item 
    \item 
    \item 
\end{enumerate}

\vspace{0.5cm}
\noindent\textbf{Exercício 13.}

\vspace{0.5cm}
\noindent\textbf{Exercício 14.}
\begin{enumerate}
    \renewcommand{\theenumi}{\alph{enumi}}
    \renewcommand{\labelenumi}{\theenumi)}
    \item 
    \item 
    \item 
    \item 
\end{enumerate}

\subsection{Função Derivada e Derivadas de Ordem Superior}

\vspace{0.5cm}
\noindent\textbf{Exercício 1.}
\begin{enumerate}
    \renewcommand{\theenumi}{\alph{enumi}}
    \renewcommand{\labelenumi}{\theenumi)}
    \item 
    \item 
    \item 
    \item 
    \item 
    \item 
\end{enumerate}

\vspace{0.5cm}
\noindent\textbf{Exercício 2.}
\begin{enumerate}
    \renewcommand{\theenumi}{\alph{enumi}}
    \renewcommand{\labelenumi}{\theenumi)}
    \item 
    \item 
\end{enumerate}

\vspace{0.5cm}
\noindent\textbf{Exercício 3.}
\begin{enumerate}
    \renewcommand{\theenumi}{\alph{enumi}}
    \renewcommand{\labelenumi}{\theenumi)}
    \item 
    \item 
    \item 
    \item 
\end{enumerate}

\subsection{Notações para a Derivada}

\vspace{0.5cm}
\noindent\textbf{Exercício 1.}
\begin{enumerate}
    \renewcommand{\theenumi}{\alph{enumi}}
    \renewcommand{\labelenumi}{\theenumi)}
    \item 
    \item 
    \item 
    \item 
    \item 
    \item 
    \item 
    \item 
    \item 
    \item 
    \item 
    \item 
    \item 
    \item 
    \item 
    \item 
\end{enumerate}

\vspace{0.5cm}
\noindent\textbf{Exercício 2.}
\begin{enumerate}
    \renewcommand{\theenumi}{\alph{enumi}}
    \renewcommand{\labelenumi}{\theenumi)}
    \item 
    \item 
\end{enumerate}

\vspace{0.5cm}
\noindent\textbf{Exercício 3.}

\vspace{0.5cm}
\noindent\textbf{Exercício 4.}

\vspace{0.5cm}
\noindent\textbf{Exercício 5.}

\vspace{0.5cm}
\noindent\textbf{Exercício 6.}

\vspace{0.5cm}
\noindent\textbf{Exercício 7.}

\vspace{0.5cm}
\noindent\textbf{Exercício 8.}
\begin{enumerate}
    \renewcommand{\theenumi}{\alph{enumi}}
    \renewcommand{\labelenumi}{\theenumi)}
    \item 
    \item 
    \item 
    \item 
    \item 
    \item 
\end{enumerate}

\vspace{0.5cm}
\noindent\textbf{Exercício 9.}

\vspace{0.5cm}
\noindent\textbf{Exercício 10.}

\vspace{0.5cm}
\noindent\textbf{Exercício 11.}

\vspace{0.5cm}
\noindent\textbf{Exercício 12.}

\vspace{0.5cm}
\noindent\textbf{Exercício 13.}

\vspace{0.5cm}
\noindent\textbf{Exercício 14.}

\vspace{0.5cm}
\noindent\textbf{Exercício 15.}

\vspace{0.5cm}
\noindent\textbf{Exercício 16.}
\begin{enumerate}
    \renewcommand{\theenumi}{\alph{enumi}}
    \renewcommand{\labelenumi}{\theenumi)}
    \item 
    \item 
\end{enumerate}

\subsection{Regra da Cadeia para Derivação de Função Composta}
Esta subseção não possui exercícios

\subsection{Aplicações da Regra da Cadeia}

\vspace{0.5cm}
\noindent\textbf{Exercício 1.}
\begin{enumerate}
    \renewcommand{\theenumi}{\alph{enumi}}
    \renewcommand{\labelenumi}{\theenumi)}
    \item 
    \item 
    \item 
    \item 
    \item 
    \item 
    \item 
    \item 
    \item 
    \item 
    \item 
    \item 
    \item 
    \item 
    \item 
    \item 
    \item 
    \item 
    \item 
    \item 
\end{enumerate}

\vspace{0.5cm}
\noindent\textbf{Exercício 2.}

\vspace{0.5cm}
\noindent\textbf{Exercício 3.}

\vspace{0.5cm}
\noindent\textbf{Exercício 4.}
\begin{enumerate}
    \renewcommand{\theenumi}{\alph{enumi}}
    \renewcommand{\labelenumi}{\theenumi)}
    \item 
    \item 
    \item 
    \item 
    \item 
    \item 
    \item 
    \item 
    \item 
    \item 
    \item 
    \item 
    \item 
    \item 
    \item 
    \item 
    \item 
    \item 
    \item 
    \item 
\end{enumerate}

\vspace{0.5cm}
\noindent\textbf{Exercício 5.}
\begin{enumerate}
    \renewcommand{\theenumi}{\alph{enumi}}
    \renewcommand{\labelenumi}{\theenumi)}
    \item 
    \item 
    \item 
    \item 
    \item 
    \item 
    \item 
    \item 
    \item 
    \item 
    \item 
    \item 
    \item 
    \item 
    \item 
    \item 
    \item 
    \item 
    \item 
    \item 
\end{enumerate}

\vspace{0.5cm}
\noindent\textbf{Exercício 6.}

\vspace{0.5cm}
\noindent\textbf{Exercício 7.}

\vspace{0.5cm}
\noindent\textbf{Exercício 8.}

\vspace{0.5cm}
\noindent\textbf{Exercício 9.}

\vspace{0.5cm}
\noindent\textbf{Exercício 10.}

\vspace{0.5cm}
\noindent\textbf{Exercício 11.}

\vspace{0.5cm}
\noindent\textbf{Exercício 12.}

\vspace{0.5cm}
\noindent\textbf{Exercício 13.}

\vspace{0.5cm}
\noindent\textbf{Exercício 14.}

\vspace{0.5cm}
\noindent\textbf{Exercício 15.}
\begin{enumerate}
    \renewcommand{\theenumi}{\alph{enumi}}
    \renewcommand{\labelenumi}{\theenumi)}
    \item 
    \item 
    \item 
    \item 
\end{enumerate}

\vspace{0.5cm}
\noindent\textbf{Exercício 16.}
\begin{enumerate}
    \renewcommand{\theenumi}{\alph{enumi}}
    \renewcommand{\labelenumi}{\theenumi)}
    \item 
    \item 
    \item 
    \item 
    \item 
    \item 
    \item 
    \item 
    \item 
    \item 
    \item 
    \item 
\end{enumerate}

\vspace{0.5cm}
\noindent\textbf{Exercício 17.}

\vspace{0.5cm}
\noindent\textbf{Exercício 18.}

\vspace{0.5cm}
\noindent\textbf{Exercício 19.}

\vspace{0.5cm}
\noindent\textbf{Exercício 20.}

\vspace{0.5cm}
\noindent\textbf{Exercício 21.}

\vspace{0.5cm}
\noindent\textbf{Exercício 22.}

\vspace{0.5cm}
\noindent\textbf{Exercício 23.}
\begin{enumerate}
    \renewcommand{\theenumi}{\alph{enumi}}
    \renewcommand{\labelenumi}{\theenumi)}
    \item 
    \item 
    \item 
\end{enumerate}

\vspace{0.5cm}
\noindent\textbf{Exercício 24.}

\vspace{0.5cm}
\noindent\textbf{Exercício 25.}

\vspace{0.5cm}
\noindent\textbf{Exercício 26.}

\vspace{0.5cm}
\noindent\textbf{Exercício 27.}

\vspace{0.5cm}
\noindent\textbf{Exercício 28.}

\vspace{0.5cm}
\noindent\textbf{Exercício 29.}
\begin{enumerate}
    \renewcommand{\theenumi}{\alph{enumi}}
    \renewcommand{\labelenumi}{\theenumi)}
    \item 
    \item 
\end{enumerate}

\vspace{0.5cm}
\noindent\textbf{Exercício 30.}

\subsection{Derivada de $f(x)^{g(x)}$}

\vspace{0.5cm}
\noindent\textbf{Exercício 1.}
\begin{enumerate}
    \renewcommand{\theenumi}{\alph{enumi}}
    \renewcommand{\labelenumi}{\theenumi)}
    \item 
    \item 
    \item 
    \item 
    \item 
    \item 
    \item 
    \item 
    \item 
    \item 
    \item 
    \item 
    \item 
    \item 
    \item 
    \item 
\end{enumerate}

\vspace{0.5cm}
\noindent\textbf{Exercício 2.}

\vspace{0.5cm}
\noindent\textbf{Exercício 3.}
\begin{enumerate}
    \renewcommand{\theenumi}{\alph{enumi}}
    \renewcommand{\labelenumi}{\theenumi)}
    \item 
    \item 
    \item 
    \item 
    \item 
    \item 
\end{enumerate}

\subsection{Derivação da Função Dada Implicitamente}

\vspace{0.5cm}
\noindent\textbf{Exercício 1.}

\vspace{0.5cm}
\noindent\textbf{Exercício 2.}

\vspace{0.5cm}
\noindent\textbf{Exercício 3.}

\vspace{0.5cm}
\noindent\textbf{Exercício 4.}
\begin{enumerate}
    \renewcommand{\theenumi}{\alph{enumi}}
    \renewcommand{\labelenumi}{\theenumi)}
    \item 
    \item 
    \item 
    \item 
    \item 
    \item 
    \item 
    \item 
    \item 
    \item 
    \item 
    \item 
\end{enumerate}

\vspace{0.5cm}
\noindent\textbf{Exercício 5.}

\vspace{0.5cm}
\noindent\textbf{Exercício 6.}

\vspace{0.5cm}
\noindent\textbf{Exercício 7.}

\vspace{0.5cm}
\noindent\textbf{Exercício 8.}
\begin{enumerate}
    \renewcommand{\theenumi}{\alph{enumi}}
    \renewcommand{\labelenumi}{\theenumi)}
    \item 
    \item 
\end{enumerate}

\vspace{0.5cm}
\noindent\textbf{Exercício 9.}

\vspace{0.5cm}
\noindent\textbf{Exercício 10.}

\vspace{0.5cm}
\noindent\textbf{Exercício 11.}

\subsection{Interpretação de $\frac{dy}{dx}$ como um Quociente Diferencial}

\vspace{0.5cm}
\noindent\textbf{Exercício 1.}
\begin{enumerate}
    \renewcommand{\theenumi}{\alph{enumi}}
    \renewcommand{\labelenumi}{\theenumi)}
    \item 
    \item 
    \item 
    \item 
\end{enumerate}

\vspace{0.5cm}
\noindent\textbf{Exercício 2.}
\begin{enumerate}
    \renewcommand{\theenumi}{\alph{enumi}}
    \renewcommand{\labelenumi}{\theenumi)}
    \item 
    \item 
\end{enumerate}

\vspace{0.5cm}
\noindent\textbf{Exercício 3.}
\begin{enumerate}
    \renewcommand{\theenumi}{\alph{enumi}}
    \renewcommand{\labelenumi}{\theenumi)}
    \item 
    \item 
\end{enumerate}

\vspace{0.5cm}
\noindent\textbf{Exercício 4.}
\begin{enumerate}
    \renewcommand{\theenumi}{\alph{enumi}}
    \renewcommand{\labelenumi}{\theenumi)}
    \item 
    \item 
\end{enumerate}

\subsection{Velocidade e Aceleração. Taxa de Variação}

\vspace{0.5cm}
\noindent\textbf{Exercício 1.}
\begin{enumerate}
    \renewcommand{\theenumi}{\alph{enumi}}
    \renewcommand{\labelenumi}{\theenumi)}
    \item 
    \item 
    \item 
    \item 
\end{enumerate}

\vspace{0.5cm}
\noindent\textbf{Exercício 2.}
\begin{enumerate}
    \renewcommand{\theenumi}{\alph{enumi}}
    \renewcommand{\labelenumi}{\theenumi)}
    \item 
    \item 
    \item 
\end{enumerate}

\vspace{0.5cm}
\noindent\textbf{Exercício 3.}
\begin{enumerate}
    \renewcommand{\theenumi}{\alph{enumi}}
    \renewcommand{\labelenumi}{\theenumi)}
    \item 
    \item 
    \item 
    \item 
\end{enumerate}

\vspace{0.5cm}
\noindent\textbf{Exercício 4.}

\vspace{0.5cm}
\noindent\textbf{Exercício 5.}

\vspace{0.5cm}
\noindent\textbf{Exercício 6.}

\vspace{0.5cm}
\noindent\textbf{Exercício 7.}
\begin{enumerate}
    \renewcommand{\theenumi}{\alph{enumi}}
    \renewcommand{\labelenumi}{\theenumi)}
    \item 
    \item 
    \item 
    \item 
\end{enumerate}

\vspace{0.5cm}
\noindent\textbf{Exercício 8.}
\begin{enumerate}
    \renewcommand{\theenumi}{\alph{enumi}}
    \renewcommand{\labelenumi}{\theenumi)}
    \item 
    \item 
    \item 
    \item 
\end{enumerate}

\vspace{0.5cm}
\noindent\textbf{Exercício 9.}
\begin{enumerate}
    \renewcommand{\theenumi}{\alph{enumi}}
    \renewcommand{\labelenumi}{\theenumi)}
    \item 
    \item 
    \item 
    \item 
    \item 
    \item 
\end{enumerate}

\vspace{0.5cm}
\noindent\textbf{Exercício 10.}
\begin{enumerate}
    \renewcommand{\theenumi}{\alph{enumi}}
    \renewcommand{\labelenumi}{\theenumi)}
    \item 
    \item 
    \item 
    \item 
\end{enumerate}

\vspace{0.5cm}
\noindent\textbf{Exercício 11.}

\vspace{0.5cm}
\noindent\textbf{Exercício 12.}

\vspace{0.5cm}
\noindent\textbf{Exercício 13.}

\vspace{0.5cm}
\noindent\textbf{Exercício 14.}
\begin{enumerate}
    \renewcommand{\theenumi}{\alph{enumi}}
    \renewcommand{\labelenumi}{\theenumi)}
    \item 
    \item 
\end{enumerate}

\vspace{0.5cm}
\noindent\textbf{Exercício 15.}

\vspace{0.5cm}
\noindent\textbf{Exercício 16.}

\vspace{0.5cm}
\noindent\textbf{Exercício 17.}

\vspace{0.5cm}
\noindent\textbf{Exercício 18.}

\vspace{0.5cm}
\noindent\textbf{Exercício 19.}

\vspace{0.5cm}
\noindent\textbf{Exercício 20.}

\vspace{0.5cm}
\noindent\textbf{Exercício 21.}

\subsection{Problemas Reta Tangente e Normal ao Gráfico de uma Função}

\vspace{0.5cm}
\noindent\textbf{Exercício 1.}
\begin{enumerate}
    \renewcommand{\theenumi}{\alph{enumi}}
    \renewcommand{\labelenumi}{\theenumi)}
    \item 
    \item 
    \item 
    \item 
\end{enumerate}

\vspace{0.5cm}
\noindent\textbf{Exercício 2.}

\vspace{0.5cm}
\noindent\textbf{Exercício 3.}

\vspace{0.5cm}
\noindent\textbf{Exercício 4.}

\vspace{0.5cm}
\noindent\textbf{Exercício 5.}

\vspace{0.5cm}
\noindent\textbf{Exercício 6.}
\begin{enumerate}
    \renewcommand{\theenumi}{\alph{enumi}}
    \renewcommand{\labelenumi}{\theenumi)}
    \item 
    \item 
\end{enumerate}

\vspace{0.5cm}
\noindent\textbf{Exercício 7.}

\vspace{0.5cm}
\noindent\textbf{Exercício 8.}

\vspace{0.5cm}
\noindent\textbf{Exercício 9.}

\vspace{0.5cm}
\noindent\textbf{Exercício 10.}

\vspace{0.5cm}
\noindent\textbf{Exercício 11.}

\vspace{0.5cm}
\noindent\textbf{Exercício 12.}

\vspace{0.5cm}
\noindent\textbf{Exercício 13.}

\vspace{0.5cm}
\noindent\textbf{Exercício 14.}

\vspace{0.5cm}
\noindent\textbf{Exercício 15.}

\subsection{Exercícios do Capítulo}

\vspace{0.5cm}
\noindent\textbf{Exercício 1.}
\begin{enumerate}
    \renewcommand{\theenumi}{\alph{enumi}}
    \renewcommand{\labelenumi}{\theenumi)}
    \item 
    \item 
    \item 
    \item 
    \item 
    \item 
    \item 
    \item 
    \item 
\end{enumerate}

\vspace{0.5cm}
\noindent\textbf{Exercício 2.}
\begin{enumerate}
    \renewcommand{\theenumi}{\alph{enumi}}
    \renewcommand{\labelenumi}{\theenumi)}
    \item 
    \item 
    \item 
    \item 
    \item 
    \item 
    \item 
    \item 
    \item 
    \item 
    \item 
    \item 
    \item 
    \item 
    \item 
    \item 
    \item 
    \item 
    \item 
    \item 
\end{enumerate}

\vspace{0.5cm}
\noindent\textbf{Exercício 3.}
\begin{enumerate}
    \renewcommand{\theenumi}{\alph{enumi}}
    \renewcommand{\labelenumi}{\theenumi)}
    \item 
    \item 
    \item 
    \item 
\end{enumerate}

\vspace{0.5cm}
\noindent\textbf{Exercício 4.}

\vspace{0.5cm}
\noindent\textbf{Exercício 5.}

\vspace{0.5cm}
\noindent\textbf{Exercício 6.}

\vspace{0.5cm}
\noindent\textbf{Exercício 7.}

\vspace{0.5cm}
\noindent\textbf{Exercício 8.}

\vspace{0.5cm}
\noindent\textbf{Exercício 9.}

\vspace{0.5cm}
\noindent\textbf{Exercício 10.}

\vspace{0.5cm}
\noindent\textbf{Exercício 11.}

\vspace{0.5cm}
\noindent\textbf{Exercício 12.}

\vspace{0.5cm}
\noindent\textbf{Exercício 13.}

\vspace{0.5cm}
\noindent\textbf{Exercício 14.}

\vspace{0.5cm}
\noindent\textbf{Exercício 15.}

\vspace{0.5cm}
\noindent\textbf{Exercício 16.}

\vspace{0.5cm}
\noindent\textbf{Exercício 17.}

\vspace{0.5cm}
\noindent\textbf{Exercício 18.}

\vspace{0.5cm}
\noindent\textbf{Exercício 19.}

\vspace{0.5cm}
\noindent\textbf{Exercício 20.}
\begin{enumerate}
    \renewcommand{\theenumi}{\alph{enumi}}
    \renewcommand{\labelenumi}{\theenumi)}
    \item 
    \item 
\end{enumerate}

\vspace{0.5cm}
\noindent\textbf{Exercício 21.}
\begin{enumerate}
    \renewcommand{\theenumi}{\alph{enumi}}
    \renewcommand{\labelenumi}{\theenumi)}
    \item 
    \item 
\end{enumerate}

\vspace{0.5cm}
\noindent\textbf{Exercício 22.}
\begin{enumerate}
    \renewcommand{\theenumi}{\alph{enumi}}
    \renewcommand{\labelenumi}{\theenumi)}
    \item 
    \item 
    \item 
    \item 
\end{enumerate}

\vspace{0.5cm}
\noindent\textbf{Exercício 23.}

\vspace{0.5cm}
\noindent\textbf{Exercício 24.}

\vspace{0.5cm}
\noindent\textbf{Exercício 25.}
\begin{enumerate}
    \renewcommand{\theenumi}{\alph{enumi}}
    \renewcommand{\labelenumi}{\theenumi)}
    \item 
    \item 
\end{enumerate}

\vspace{0.5cm}
\noindent\textbf{Exercício 26.}
\begin{enumerate}
    \renewcommand{\theenumi}{\alph{enumi}}
    \renewcommand{\labelenumi}{\theenumi)}
    \item 
    \item 
\end{enumerate}

\vspace{0.5cm}
\noindent\textbf{Exercício 27.}
\begin{enumerate}
    \renewcommand{\theenumi}{\alph{enumi}}
    \renewcommand{\labelenumi}{\theenumi)}
    \item 
    \item 
\end{enumerate}

\vspace{0.5cm}
\noindent\textbf{Exercício 28.}
\begin{enumerate}
    \renewcommand{\theenumi}{\alph{enumi}}
    \renewcommand{\labelenumi}{\theenumi)}
    \item 
    \item 
\end{enumerate}

\vspace{0.5cm}
\noindent\textbf{Exercício 29.}
\begin{enumerate}
    \renewcommand{\theenumi}{\alph{enumi}}
    \renewcommand{\labelenumi}{\theenumi)}
    \item 
    \item 
\end{enumerate}

\vspace{0.5cm}
\noindent\textbf{Exercício 30.}

\vspace{0.5cm}
\noindent\textbf{Exercício 31.}

\vspace{0.5cm}
\noindent\textbf{Exercício 32.}

\vspace{0.5cm}
\noindent\textbf{Exercício 33.}
\begin{enumerate}
    \renewcommand{\theenumi}{\alph{enumi}}
    \renewcommand{\labelenumi}{\theenumi)}
    \item 
    \item 
\end{enumerate}

\vspace{0.5cm}
\noindent\textbf{Exercício 34.}

\vspace{0.5cm}
\noindent\textbf{Exercício 35.}

\vspace{0.5cm}
\noindent\textbf{Exercício 36.}
\begin{enumerate}
    \renewcommand{\theenumi}{\alph{enumi}}
    \renewcommand{\labelenumi}{\theenumi)}
    \item 
    \item 
\end{enumerate}

\vspace{0.5cm}
\noindent\textbf{Exercício 37.}

\vspace{0.5cm}
\noindent\textbf{Exercício 38.}

\vspace{0.5cm}
\noindent\textbf{Exercício 39.}
\begin{enumerate}
    \renewcommand{\theenumi}{\alph{enumi}}
    \renewcommand{\labelenumi}{\theenumi)}
    \item 
    \item 
    \item 
    \item 
\end{enumerate}

\vspace{0.5cm}
\noindent\textbf{Exercício 40.}

\vspace{0.5cm}
\noindent\textbf{Exercício 41.}
\begin{enumerate}
    \renewcommand{\theenumi}{\alph{enumi}}
    \renewcommand{\labelenumi}{\theenumi)}
    \item 
    \item 
    \item 
    \item 
    \item 
    \item 
\end{enumerate}

\vspace{0.5cm}
\noindent\textbf{Exercício 42.}

\vspace{0.5cm}
\noindent\textbf{Exercício 43.}

\vspace{0.5cm}
\noindent\textbf{Exercício 44.}

\vspace{0.5cm}
\noindent\textbf{Exercício 45.}

\vspace{0.5cm}
\noindent\textbf{Exercício 46.}

\vspace{0.5cm}
\noindent\textbf{Exercício 47.}

\vspace{0.5cm}
\noindent\textbf{Exercício 48.}

\vspace{0.5cm}
\noindent\textbf{Exercício 49.}

\end{document}
